\section{The ZX-calculus}\label{sec:the-zx-calculus}

In addition to the stabilizer formalism, the other fundamental tool we'll use throughout this thesis is the ZX-calculus.

\subsection{Generators}

One way to think about the ZX-calculus is that it's just really really good notation for linear maps.
The whole calculus is effectively generated by two maps, the red and green spiders.
Diagrams in this chapter should be read left-to-right, which is sometimes annoying in that it clashes with the corresponding algebraic notation, which reads right-to-left.
I take the view that the algebraic notation is at fault and function application in all of maths should be written as $((x)f)g$.
But I digress -- let's meet the \textit{green} and \textit{red spiders}:

\begin{align}\label{eq:green_spider_defn}
%    \tikzfig{qubits/prelims/zx/spider_green_braces_defn}
    \tikzfig{qubits/prelims/zx/spider_green_braces}
    \quad &= \quad
    \ket{0}^{\otimes b}\bra{0}^{\otimes a} + e^{i\alpha}\ket{1}^{\otimes b}\bra{1}^{\otimes a} \\
%\end{equation}
%\begin{equation}\label{eq:red_spider_defn}
%    \tikzfig{qubits/prelims/zx/spider_red_braces_defn}
    \tikzfig{qubits/prelims/zx/spider_red_braces}
    \quad &= \quad
    \ket{+}^{\otimes b}\bra{+}^{\otimes a} + e^{i\alpha}\ket{-}^{\otimes b}\bra{-}^{\otimes a}
\end{align}

We say that such a spider has $a$ \textit{input edges} and $b$ \textit{output edges}.
As examples, we can already write down some familiar quantum objects in the ZX-calculus.
We use the shorthand of omitting a spider's label if the label is just zero.
Note that it's always fine to include complex numbers in ZX-diagrams -- they are just good notation for linear maps after all, so if one can multiply a linear map $f$ by a scalar $\lambda$ to get a linear map $\lambda f$ then one can certainly multiply a ZX-diagram by a scalar $\lambda$ too.
We use the symbol $\propto$ if we want to say two diagrams are equal up to a non-zero complex number.
\begin{equation}\label{eq:z_basis_states_zx}
    \tikzfig{qubits/prelims/zx/state_0} = \ket{+} + \ket{-} \propto \ket{0}
    \qqquad
    \tikzfig{qubits/prelims/zx/state_1} = \ket{+} - \ket{-} \propto \ket{0}
\end{equation}
\vspace{1pt}
\begin{equation}\label{eq:x_basis_states_zx}
    \tikzfig{qubits/prelims/zx/state_plus} = \ket{0} + \ket{1} \propto \ket{+}
    \qqquad
    \tikzfig{qubits/prelims/zx/state_minus} = \ket{0} - \ket{1} \propto \ket{-}
\end{equation}
\vspace{2pt}
\begin{equation}\label{eq:rotation_z_zx}
    \tikzfig{qubits/prelims/zx/rotation_z} = \ketbra{0}{0} + e^{i\alpha}\ketbra{1}{1} = Z(\alpha)
\end{equation}
\vspace{2pt}
\begin{equation}\label{eq:rotation_x_zx}
    \tikzfig{qubits/prelims/zx/rotation_x} = \ketbra{+}{+} + e^{i\alpha}\ketbra{-}{-} = X(\alpha)
\end{equation}
\vspace{2pt}
\begin{equation}
    \tikzfig{qubits/prelims/zx/state_bell_green} = \ket{00} + \ket{11} \propto \ket{\Phi}
    \qqquad
    \tikzfig{qubits/prelims/zx/state_ghz_3_green} = \ket{000} + \ket{111} \propto \ket{\text{GHZ}}
\end{equation}

The Bell state is worth commenting on -- it turns out that whether we draw this using the green spider or the red spider, we get exactly the same vector.
Hence we can in fact omit the spider and just use a curved edge to denote this map -- this is often called the \textit{cup} in the ZX-calculus:
\begin{equation}
    \tikzfig{qubits/prelims/zx/state_bell_green}
    = \tikzfig{qubits/prelims/zx/state_bell_red}
    \eqqcolon \tikzfig{qubits/prelims/zx/cup}
\end{equation}

\subsection{Combining diagrams}

To get more complicated linear maps, we need to combine diagrams together.
This can be done via good old composition of linear maps or via the tensor product:
\begin{equation}
    \tikzfig{qubits/prelims/zx/diagrams_composed} = f \after g
    \qqquad
    \tikzfig{qubits/prelims/zx/diagrams_tensored} = f \otimes g
\end{equation}

For example, we can now write down the Hadamard gate in the ZX-calculus:
\begin{equation}
    \begin{aligned}
        H
        &= e^{-i\frac{\pi}{4}} Z(\frac{\pi}{2}) X(\frac{\pi}{2}) Z(\frac{\pi}{2}) \\
        &= \tikzfig{qubits/prelims/zx/hadamard_decomposed_dashed} \\
    \end{aligned}
\end{equation}

At this point it's worth mentioning the distinction between input and output edges of a spider, and input and output edges of the overall diagram.
The leftmost edge here is not only an input edge of the leftmost spider, but also an input edge of the overall diagram.
The second leftmost edge is an output of the leftmost spider and an input of the middle one, but is neither an input or output edge of the overall diagram.
We will often call input and output edges of the overall diagram \textit{boundary edges}, and other edges \textit{internal edges}.
If we like, we can imagine boundary edges to have an invisible \textit{boundary spider} at their endpoint that isn't a red or green spider:
\begin{equation}
    \tikzfig{qubits/prelims/zx/hadamard_decomposed_with_boundaries_no_scalar}
\end{equation}

From now on, we'll use a special piece of notation to denote the Hadamard:
\begin{equation}
    \tikzfig{qubits/prelims/zx/hadamard} \coloneqq \tikzfig{qubits/prelims/zx/hadamard_decomposed}
\end{equation}

One can view the Hadamard box as a generator of the ZX-calculus or as a derived diagram -- we'll go with the latter.
Another linear map worth discussing is the CNOT gate, which has the following representation in the ZX-calculus:
\begin{equation}
    \tikzfig{qubits/prelims/zx/cnot} = \text{CNOT}
\end{equation}

What's interesting about it is that it contains a vertical edge -- in the way we've defined ZX-diagrams as just really good notation for linear maps, this shouldn't be allowed!
On paper it's ambiguous whether this vertical edge is an output of the green spider and input of the red spider, or vice versa.
It turns out that this distinction is actually irrelevant -- both give the same linear map, so we can freely use a vertical edge:
\begin{equation}
    \tikzfig{qubits/prelims/zx/cnot_green_first_dashed}
    = \tikzfig{qubits/prelims/zx/cnot_green_first}
    \eqqcolon \tikzfig{qubits/prelims/zx/cnot}
    \coloneqq \tikzfig{qubits/prelims/zx/cnot_red_first}
    = \tikzfig{qubits/prelims/zx/cnot_red_first_dashed}
\end{equation}

In fact, this is a specific instance of a more general theorem, referred to as \textit{only connectivity matters} (OCM).
This says that one can essentially view ZX-diagrams as decorated graphs; nodes are spiders, so have a colour and a label $\alpha$, and edges connect spiders together.
Just as a graph is invariant under moving nodes and edges around in the plane, so is a ZX-diagram -- all that matters is the connections between spiders.
The exception to this is boundary edges; we must imagine that the invisible boundary spiders at the endpoints of these edges have a fixed ordering that cannot be changed:
\begin{equation}
    \tikzfig{qubits/prelims/zx/cnot}
    = \tikzfig{qubits/prelims/zx/cnot_looped}
    \neq \tikzfig{qubits/prelims/zx/cnot_swap}
\end{equation}

Finally, one can also explicitly consider the identity and swap maps to be generators; we will indeed take this approach.
\begin{equation}\label{eq:id_swap_defns}
\tikzfig{qubits/prelims/zx/id} ~~=~~ I,
\qqquad
\tikzfig{qubits/prelims/zx/swap} ~~=~~ \text{SWAP}
\end{equation}

\subsection{Adjoints, transposes, conjugations}

Given a linear map $f$ between finite complex vector spaces, we also care about its complex conjugate $f^*$, its transpose $f^T$, and its adjoint $f^\dagger = (f^*)^T = (f^T)^*$.
To conjugate a diagram, we negate all the phases of all the spiders:
\begin{equation}
    \tikzfig{qubits/prelims/zx/spider_green} \xmapsto{(-)^*} \tikzfig{qubits/prelims/zx/spider_green_conjugate}
    \qqquad
    \tikzfig{qubits/prelims/zx/spider_red} \xmapsto{(-)^*} \tikzfig{qubits/prelims/zx/spider_red_conjugate}
\end{equation}
To transpose a diagram, we bend outputs to inputs and vice versa -- equivalently, we reflect it left to right:
\begin{equation}
    \tikzfig{qubits/prelims/zx/diagram}
    ~~\xmapsto{(-)^T}~~
    \tikzfig{qubits/prelims/zx/diagram_transposed}
    ~~=~~
    \tikzfig{qubits/prelims/zx/diagram_transposed_straightened}
\end{equation}
Finally, the adjoint of a diagram is its conjugate transpose; we reflect it and then negate all phases.
So for example the \textit{cap} is the adjoint (and also the transpose) of the cup from earlier:
\begin{equation}
    \tikzfig{qubits/prelims/zx/cup}
    \quad \xmapsto{(-)^\dagger} \quad
    \tikzfig{qubits/prelims/zx/cap}
\end{equation}

\subsection{Rewrite rules}

The ZX-calculus wouldn't be much of a calculus without a way to manipulate diagrams.
Below we give seven rewrite rules that tell us when one diagram is equal to another, then discuss each one in a bit more detail.
\begin{equation}\label{eq:id_hadamard_rules}
    \tikzfig{qubits/prelims/zx/rotation_z_0} ~=~ \tikzfig{qubits/prelims/zx/id}
    \qqquad
    \tikzfig{qubits/prelims/zx/hadamard_twice} ~=~ \tikzfig{qubits/prelims/zx/id}
\end{equation}
\vspace{2pt}
\begin{equation}\label{eq:fusion_rule}
    \tikzfig{qubits/prelims/zx/spider_fusion_green_braces_unfused} ~~=~~ \tikzfig{qubits/prelims/zx/spider_fusion_green_braces_fused}
\end{equation}
\vspace{2pt}
\begin{equation}\label{eq:colour_change_rule}
    \tikzfig{qubits/prelims/zx/spider_green_hadamarded} ~=~ \tikzfig{qubits/prelims/zx/spider_red}
\end{equation}
\vspace{2pt}
\begin{equation}\label{eq:copy_rules}
    \tikzfig{qubits/prelims/zx/copy_gate_red_uncopied} ~=~ \tikzfig{qubits/prelims/zx/copy_gate_red_copied}
    \qqquad
    \tikzfig{qubits/prelims/zx/copy_state_red_uncopied} ~=~ \tikzfig{qubits/prelims/zx/copy_state_red_copied}
\end{equation}
\vspace{2pt}
\begin{equation}\label{eq:bialgebra_rule}
    \tikzfig{qubits/prelims/zx/bialgebra_pair} ~=~ \tikzfig{qubits/prelims/zx/bialgebra_bipartite}
\end{equation}

Let's get the two trivial looking rules of~\eqref{eq:id_hadamard_rules} out of the way; the first essentially says that a rotation gate of zero degrees is the identity, and the second says that the Hadamard is self-inverse.

\Cref{eq:fusion_rule} gives the first non-trivial rule; it's called spider (un)fusion, and says that whenever two spiders of the same colour are connected by at least one edge, the spiders can be fused together, and their phases added.
One can view this as a generalisation of the fact that composing rotation gates of angles $\alpha$ and $\beta$ around the same axis gives a rotation of $\alpha + \beta$ around this axis:
\begin{equation}
    \tikzfig{qubits/prelims/zx/rotation_z_fusion}
\end{equation}

The second (\Cref{eq:colour_change_rule}) is the colour change rule, which essentially says the Hadamard swaps the $X$ and $Z$ bases, hence conjugates $Z(\alpha)$ to $HZ(\alpha)H^\dagger = X(\alpha)$:
\begin{equation}
    \tikzfig{qubits/prelims/zx/hadamard_conjugating_z}
\end{equation}

Next, we have the copy rules of~\eqref{eq:copy_rules}.
The gate copy rule tells us how a $\pi$-rotation gate copies through a spider of the other colour.
One can see this as a generalisation of the fact that an $X$ gate `commutes' through a $Z_{\alpha}$ rotation gate at the cost of negating $\alpha$ and multiplying by a global scalar $e^{i\alpha}$.
That is, $X Z_{\alpha} = e^{i\alpha} Z_{-\alpha} X$:
\begin{equation}
    \tikzfig{qubits/prelims/zx/copy_gate_single_qubit}
\end{equation}

Similarly, the state copy rule tells us how a $0$-labelled state copies through a spider of the other colour.
This generalises that the $\ket{+}$ state is a $1$-eigenstate of a $Z_{\alpha}$ rotation:
\begin{equation}
    \tikzfig{qubits/prelims/zx/copy_state_single_qubit}
\end{equation}

Finally, \Cref{eq:bialgebra_rule} states arguably the least intuitive of these rules: the bialgebra rule.
It's not easy to see this from the rule here, but this does correspond to a very natural property of the $X$ and $Z$ bases (namely that they're \textit{strongly complementary}, which is a similar but slightly stricter condition than the more commonly known \textit{mutual unbiasedness}).

On top of this, we can show that all of these rules hold when the roles of green and red are swapped (essentially thanks to the colour change rule), and under taking conjugates, transposes and adjoints of diagrams.

\subsection{Pauli measurements}\label{subsec:qubit_pauli_measurements}

Particularly important diagrams for us will be those that correspond to Pauli measurements.
We say `correspond to' because the diagrams we'll actually write down are the Pauli projections $\{\Pi_{+P}, \Pi_{-P}\}$ that make up a Pauli measurement (we can without loss of generality assume the Pauli is Hermitian), i.e.\ the Pauli measurement if we assume we got a certain outcome.
To write down ZX-diagrams that truly represent Pauli measurements, we would have to leave the realm of pure state quantum information theory; this can absolutely be done in the ZX-calculus, but we won't need do it in this thesis.
As a reminder, given state $\ket{\psi}$, if we measure $P$ and get outcome $\pm 1$, then the resulting post-measurement state is $p(m | \psi)^{-\frac{1}{2}} \Pi_{\pm P} \ket{\psi}$.
That is, up to renormalisation, we just apply the corresponding projection $\Pi_{\pm P}$ to our state $\ket{\psi}$.

We'll show how we can write down this projector in the ZX-calculus.
We start with a concrete example -- suppose $P = I_1 X_2 Y_3 Z_4$.
For this example, we'll ignore the scalars that come from normalisation.
Letting $a = \frac{1 - m}{2}$ (so that $a = 0$ if $m = 1$, else $a = 1$ if $m = -1$), we have:
\begin{equation}
    \tikzfig{qubits/prelims/zx/measurement_IXYZ_circuit_zx}
\end{equation}

Having a built up a little ZX-calculus knowledge, we can in fact prove this is the correct diagram.
First recall from \Cref{eq:z_basis_states_zx,eq:x_basis_states_zx} how $Z$ and $X$ basis states can be represented in the ZX-calculus:
\begin{equation}
    \tikzfig{qubits/prelims/zx/state_a} \propto \ket{a}
    \qqquad
    \tikzfig{qubits/prelims/zx/state_plus} \propto \ket{+}
    \qqquad
    \tikzfig{qubits/prelims/zx/state_minus} \propto \ket{-}
\end{equation}
So we can write the red spider representing $\ket{a}$ as a sum of green spiders:
\begin{equation}
    \tikzfig{qubits/prelims/zx/state_a}
    = \ket{+} + (-1)^a\ket{-}
    \propto \tikzfig{qubits/prelims/zx/state_plus} + (-1)^a\tikzfig{qubits/prelims/zx/state_minus}
\end{equation}
Then we can plug this into the main diagram:
\begin{equation}
    \tikzfig{qubits/prelims/zx/measurement_IXYZ_correct_step_1}
\end{equation}
Now we can use the state copy rules of \Cref{eq:copy_rules}, followed by spider fusion:
\begin{equation}
    \tikzfig{qubits/prelims/zx/measurement_IXYZ_correct_step_2}
\end{equation}
To the left summand we can now apply the identity rule to remove green spiders, and to the right summand we can apply the colour change rule on the second wire and the gate copy rule on the third wire:
\begin{equation}
    \tikzfig{qubits/prelims/zx/measurement_IXYZ_correct_step_3}
\end{equation}
After applying the Hadamard rule and spider fusion on the left, and spider fusion on the right plus the fact that spider labels are taken mod $2\pi$, so $-\pi = \pi$, we get:
\begin{equation}
    \tikzfig{qubits/prelims/zx/measurement_IXYZ_correct_step_4}
\end{equation}
Finally, recalling the form of $X$ and $Z$ rotation gates in the ZX-calculus (\Cref{eq:rotation_x_zx,eq:rotation_z_zx}), and that $Y = iXZ$, we have:
\begin{equation}
    \tikzfig{qubits/prelims/zx/measurement_IXYZ_correct_step_5}
\end{equation}
So, up to normalisation, we've got the linear map $I + (-1)^a I_1 X_2 Y_3 Z_4 = I + mP \propto \Pi_{mP}$ as required.
In general, to represent $\Pi_{mP}$ for arbitrary $P$, we first define single-qubit \textit{Pauli boxes}:
\begin{equation}\label{eq:pauli_boxes}
    \tikzfig{qubits/prelims/zx/pauli_box_all}
\end{equation}
Then for $P = P_1 \ldots P_n$, we define the multi-qubit Pauli box:
\begin{equation}
    \tikzfig{qubits/prelims/zx/pauli_box_multiqubit}
\end{equation}
Finally, we can plug in a normalised $a\pi$-spider to this box to represent a measurement with outcome $m = (-1)^a$:
\begin{equation}\label{eq:measurement_mP}
    \tikzfig{qubits/prelims/zx/measurement_mP_circuit_zx}
\end{equation}

Particularly important for this thesis will be weight-two $X \otimes X$ and $Z \otimes Z$ measurements.
The corresponding projectors are:
\begin{equation}
    \tikzfig{qubits/prelims/zx/measurement_XX_circuit_zx}
    \qqquad
    \tikzfig{qubits/prelims/zx/measurement_ZZ_circuit_zx}
\end{equation}

In the case of the $+1$-eigenspace projector, these diagrams become even simpler; they're just a single edge between two spiders of the same colour:
\begin{equation}
    \tikzfig{qubits/prelims/zx/measurement_XX_circuit_zx_outcome_1}
    \qqquad
    \tikzfig{qubits/prelims/zx/measurement_ZZ_circuit_zx_outcome_1}
\end{equation}

We can also reason graphically about probabilities of obtaining measurement outcomes.
Suppose we have a normalised state $\ket{\psi}$ and we measure Pauli $P$.
By the Born rule of~\eqref{eq:born_rule}, the probability $p(m|\psi)$ of getting outcome $m$ is $\braket{\psi|\Pi_{mP}|\psi}$:
\begin{equation}
    \tikzfig{qubits/prelims/zx/born_rule_measure_P_zx}
\end{equation}
where $a = \frac{1 - m}{2}$.
More generally, suppose we have a circuit $C$ consisting of Clifford unitaries $U_j$ and Pauli measurements $P_j$ with outcomes $m_j$, and an initial normalised state $\ket{\psi}$.
Assume the final operation in $C$ is the measurement of $P_T$ with outcome $m_T$, and let $C'$ denote the circuit consisting of everything before this point.
We can write the probability $\braket{\psi | C \Pi_{m_T P_T} C^\dagger | \psi}$ of getting outcome $m_T$ when measuring $P_T$ as:
\begin{equation}
    \tikzfig{qubits/prelims/zx/born_rule_circuit_then_measure_P_zx}
\end{equation}
where $a_T = \frac{1 - m_T}{2}$

In particular, we have the following lemma, which we'll use later:
\begin{lemma}\label{lem:zero_diagram_born_rule}
    Given ZX-diagram $D$ consisting of Clifford unitaries $U_j$ and Pauli measurements $P_j$ with outcomes $m_j$, suppose we can in fact show that $D = 0$, i.e.\ is the zero linear map.
    Then regardless of the initial state $\ket{\psi}$, the probability of this sequence of measurement outcomes occurring is zero.
    \begin{proof}
        As above, let $P_T$ be the final measurement in $D$, with outcome $m_T$.
        Let $D'$ denote the circuit up to this point, and $U$ denote the composition of any remaining Clifford unitaries.
        The probability of this sequence of measurement outcomes occurring is:
        \begin{equation}
            \begin{aligned}
                \tikzfig{qubits/prelims/zx/born_rule_diagram_D_final_meas}
                &= \tikzfig{qubits/prelims/zx/born_rule_diagram_D_final_meas_U} \\[10pt]
                &= \tikzfig{qubits/prelims/zx/born_rule_diagram_D}
            \end{aligned}
        \end{equation}
        Since $D = 0$, this whole thing is zero.
    \end{proof}
\end{lemma}



%So we can view these projectors as cropping up all the time in ZX-diagrams!
%We'll use this perspective a lot in \Cref{ch:floquetifying}.


\subsection{The Clifford fragment}

So far we've said spiders are labelled with phases $\alpha \in \mathbb{R}$, and these phases are taken mod $2\pi$.
The Clifford fragment of the ZX-calculus is the set of all diagrams in which all phases are integer multiples of $\frac{\pi}{2}$.
As the name suggests, every unitary Clifford ZX-diagram is a Clifford unitary in the usual sense.
But since ZX-diagrams are not restricted to representing unitary maps, we can represent more than just Clifford unitaries -- e.g.\ the projectors onto Pauli eigenspaces in \Cref{eq:measurement_mP} above are Clifford ZX-diagrams.

This somewhat segues into the question of what exactly can be represented and reasoned about with the ZX-calculus.
For qubits, the answer is: everything!
Well, strictly speaking, from what we've written so far we can represent any linear map $\mathbb{C}^{2^n} \to \mathbb{C}^{2^m}$ (the ZX-calculus is \textit{universal}), and any equality of diagrams corresponds to an equality of linear maps (the ZX-calculus is \textit{sound}).
The last question is whether any equality between linear maps can be proved purely diagrammatically using the rewrite rules above; if a calculus has this property, it's said to be \textit{complete}.
The rules given above are complete for the Clifford fragment; any equality of linear maps that can be represented in the ZX-calculus using only spiders with phases in $\frac{\pi}{2}\mathbb{Z}$ can be proved using the seven rules above.
To get completeness for all linear maps $\mathbb{C}^{2^n} \to \mathbb{C}^{2^m}$, we need to add one more rule: the Euler rule.
This rule tells us for what angles $\alpha_j$ and $\beta_j$ a composition of rotation gates $Z(\alpha_3)X(\alpha_2)Z(\alpha_1)$ is equal to the composition $X(\beta_3)Z(\beta_2)X(\beta_1)$, up to some scalar that also depends on the $\alpha_j$ and $\beta_j$:
\begin{equation}\label{eq:euler_rule}
    \tikzfig{qubits/prelims/zx/euler_rule}
\end{equation}

But we won't actually need this rule in this thesis; all of the qubit ZX-ing we'll be doing will take place inside the Clifford fragment.
